\chapter{Ringraziamenti}

Desidero ringraziare mia madre, la donna che ha dato tutto per me da quando mi ha messo al mondo, l'unica che per me farebbe davvero di tutto, colei che mi ha insegnato ad avere dei valori, a non fermarmi d'avanti a niente e di non avere mai timore di quale possa essere il futuro ma di crearlo con le proprie mani. A lei che possiede una forza che non ho mai visto in nessuno, capace di mettersi in secondo piano e di sacrificarsi anche per regalarci un semplice sorriso e non farci pesare nulla, di nasconderci dalle paure e dalle incombenze che la vita ci ha riservato. Sapete mia madre ha sempre lottato, prima della nascita mia e di mio fratello, con la sua famiglia, poi con l'avvento di noi due, con la nostra di famiglia e tutt'ora alla sua età continua a combattere, a non farsi mettere i piedi in testa da nessuno e a farsi valere in ogni decisione e in ogni situazione che richiede una mente lucida e armata di santa pazienza. Non si ferma mai nemmeno quando è sfinita, stanca, distrutta dal lavoro e dai servizi che porta a termine ogni santo giorno. Ebbene è la persona che ispira il mio sogno più grande ossia diventare un uomo, e chi meglio di una donna, della propria madre può insegnare al figlio maschio a divenire tale, come sono cresciuto, come ho vissuto e come ho approcciato alla vita e come la affronto a testa alta oggi, dalle sue azioni ho capito cosa vuol dire sacrificarsi, dare il tutto per tutto sia ad un ideale, o ad una persona, il coraggio che ci vuole nel prendere una decisione ed affrontarne le conseguenze senza mai tirarsi indietro. Sembra che stia parlando di una persona capace di fare l'impossibile, beh per i suoi figli è capace di far avverare l'utopia. Sono cresciuto comunque in un ambiente che può racchiudersi in un ossimoro, sarà stato difficile crescere due figli maschi per la maggior parte da sola, dove ogni secondo di distrazione poteva essere fatale nella nostra incolumità o nel modo in cui poteva sviarci in altre strade ben poco rosee, ma, lei è stata in grado di far fronte a tutto senza mai aspettarsi nulla in cambio. Ci sono state volte in cui l'ho delusa ne sono certo, come quando andai alle medie dal doposcuola perché non volevo studiare, ma non mi ha mai fatto pesare la cosa anzi, affrontava la questione sempre con severità dandomi però ogni volta una scelta o la scuola o il lavoro e di questo le sono totalmente grato, credo che il suo approccio ossia quello di non aver mai obbligato nessuno di noi due figli a fare qualcosa ma a darci la piena libertà sia stato essenziale. Come ho detto è lei la musa ispiratrice che mi porta a non arrendermi, a dare il meglio di me, a fare ogni volta un passo in avanti anche quando la strada non è stata ancora tracciata, un uomo capace di suscitare conforto, speranza, che possa dare tranquillità alle persone amate non per far piacere a qualcuno o essere denominato come "bravo ragazzo", non per dio, non per essere denominato santo o avere l'ardire di impuntarmi nell'essere il migliore, mai fatto e mai lo farò, voglio solo che le persone intorno a me sentano ciò che io ho sentito ogni volta che vedevo mia madre affrontare la vita, una forza in grado di proteggere, questo è quello che voglio diventare, un uomo in grado di proteggere, di mettere il proprio petto a servizio come uno scudo e chi meglio di mia madre poteva insegnarmi ciò, colei che ci ha sempre protetti da qualsiasi ombra di questa vita e che ha permesso a me e a mio fratello di diventare uomini. A te mamma devo ogni mio sforzo e ogni mio traguardo per sempre.\\

A mio fratello, tra di noi non scorre sempre buon sangue, a volte abbiamo litigato pesantemente, a volte invece facciamo finta che uno di noi due non esista in casa, ci isoliamo nei nostri mondi ognuno con la consapevolezza che l'altro stia bene o che comunque l'altro non abbiamo chissà che problemi. Invece di problemi ne abbiamo eh? Na cifra. Non ho mai voluto romperti le scatole su ciò che ho passato, avevi le tue preoccupazioni e non volevo addossarti sulle tue spalle anche le mie, non so se ti ricordi ma non parlavamo mai se non quando dovevamo romperci le scatole a vicenda, in realtà sono pochissime le volte che mi sono aperto con te, non perché non volessi eh ripeto ma perché è un tratto mio non parlare, anche se ultimamente questa cosa sta mano mano scemando me la porto comunque ancora dietro. So già che o ti arrabbierai o sarà l'ennesima delusione che ti provocherò il fatto che non parlerò di papà, resterò coerente con me stesso, del resto non voglio sembrare ipocrita. Ora mi dirai che razza di ringraziamento è un discorso dove non ci siano ringraziamenti, hai ragione, con tutto quello che abbiamo passato non avrò rimpianti o desideri verso di te, sei il fratello di cui avevo bisogno, ormai mi conosci sai che ho bisogno dei miei spazi e il silenzio è l'arma che uso quando la testa esplode di pensieri, non mi hai mai assillato e questo per me è stato fondamentale, mi ha permesso di percepire questa casa come luogo sicuro per qualsiasi sfida dovessi affrontare, non mi hai mai giudicato per essere in ritardo con l'università, con la macchina, in generale con la vita (a parte la questione papà) ma del resto su questo abbiamo due visioni diverse e ovviamente rispetto la tua. So che in questo periodo stai attraversando parecchie delusioni ma sei forte, sei talmente forte che continui ad avanzare, tratto distintivo di nostra madre, a qualsiasi torto ti venga fatto. Quella forza che ho citato prima riferendomi a nostra madre la rivedo in te, per questo sei colui che io devo superare, con i miei difetti ovvio sono pur sempre una persona completamente diversa da te ma in quanto uomo è mio dovere superarti, hai ciò che io desidero: forza,coraggio,resistenza, mente matematica, anche bellezza eh, robustezza insomma incuti fiducia nelle persone, ciò che io invece non riesco e lo percepisco. E' come se al nostro posto doveva nascere un singolo individuo con le capacità di entrambi, d'altronde tu hai ciò che io non ho e io ho ciò che tu non hai tipo la pazienza nello studio, etc... sta di fatto che sei colui che non reputo solo come mio fratello maggiore ma il limite che devo superare, ecco perché quella sfida è ancora in piedi. Insomma il mio ringraziamento è quello di avere te, non mi interessa se sarai arrabbiato con me, o se mi odierai a tal punto che in un futuro prossimo non vorrai sentire la mia voce ne tanto meno vedermi in faccia, per quello che posso la mia mano sarà tesa per te, su questo devo fare un appunto però, sono belle parole ma non le ho messe sempre in pratica, in primis questione papà, non ti ho aiutato quando ne avevi bisogno, non mi scuso di questo, le mie ragioni le avevo, forse su questo fronte devo migliorare ma più vado avanti e più capisco che devo fare la mia parte, per diventare un uomo devo tendere la mano non alle persone che odio ma a coloro che non mi hanno fatto niente e che da me volevano solamente un aiuto.
A tutto ciò che hai perso e a quello che hai guadagnato non osare arrenderti, voglio superarti al meglio delle tue possibilità.\\


ai miei zii (famiglia Aversano),
mi avete dato tanto coraggio e speranza nel percorrere questa strada, ultimamente non credevo nemmeno di farcela a finire, ormai era alienante venir bocciato sempre per la stessa materia e ho pensato parecchie volte di rinunciare. Non l'ho mai detto anche perché sono pensieri di circostanza quando qualcuno subisce la stanchezza e i pensieri cominciano a prendere il sopravvento, tant'è che mi informai anche nel cambiare università per un singolo esame, menomale che non l'ho fatto. Vi ringrazio per il supporto che mi avete dato, del desiderio di vedermi realizzare questo sogno che ormai è diventato realtà, la vostra vicinanza è stata per me fonte di resilienza e controllo nell'affrontare quel maledetto ultimo esame, per questo zia Rosaria, zio Rosario, Salvatore e Martina vi ringrazio dal profondo del cuore.\\

questa parte la vorrei dedicare a due famiglie che sono entrate a far parte della mia vita (ringraziando a Dio) e che mi hanno trattato come un figlio e come un fratello.\\

Parto prima da un mio fratello non consanguineo ma quello solo manca ormai, Nikol Pellecchia.
Nikol ormai ci conosciamo da 10 anni quasi 11, ne abbiamo passate di tutti i colori, infatti meglio se non ricordo in che guai ci siamo imbarcati. Mi hai conosciuto che ero un ragazzo silenzioso, immobile, quasi impercettibile, in quel periodo come ben ricordi stavo attraversando un tornado di brutte situazioni, sia a scuola che a casa. Sei l'unico che mi ha visto nel mio momento di buio, nel fondo di un pozzo dove non vedevo la luce. Sei stato il primo che mi ha accettato anche quando tutti mi scartavano, non volevi chissà che cosa da me ma ciò che percepivo era la forma più pura di amicizia spontanea che avessi mai avuto, mi hai tirato fuori da quel pozzo e mi hai teso una mano che mai dimenticherò facendomi ringraziare per la prima volta il signore per avermi fatto incontrare una persona del genere, quando prima ne odiavo la maggior parte. Non abbiamo mai litigato, anche se tu dici 1 volta nemmeno me la ricordo, sarà stato qualche impiccio durato niente anche perché mi ricordo solo al quinto anno di superiori di aver deciso di non aspettarti più sotto casa perché svegliarmi alle 6 per aspettare 25 minuti non era proprio il massimo. Ok detta così sembra che stia raccontando di una persona che comunque mi faceva passare i guai, ed è così, ma se tornassi indietro non la cambierei mai e poi mai. Dal momento che ti ho incontrato ho capito una cosa fondamentale, a volte i detti antichi possono sbagliare, e ora mi dirai -che vuol dire? conosci molto bene il detto "fidarsi è bene non fidarsi è meglio", ca****e dico io, come faccio a non fidarmi di te che hai sempre, e sottolineo sempre, un occhio di riguardo nei miei confronti, quando la mia situazione economica non mi permetteva di fare chissà che e non volevo scendere per non pesare sulle spalle di qualcuno, quando non mi facevo sentire non perché non volessi ma già mi conosci e sai che io e whatsapp siamo due cose separate e mi tempestavi di chiamate per convincermi a scendere anche per andare a giocare una bolletta e finivamo nei peggio guai, quando mi includevi in qualsiasi cosa, come si può dico io non fidarsi di una persona che non ti ha mai voluto per quello che hai tipo soldi oppure prestigio o altro, che c'è sempre stata nei momenti del bisogno e mai si è tirata indietro, dico io come si fa bah, non ci sono parole a tutto ciò che vorrei dirti ma una cosa è certa mai e dico mai volterò le spalle a te o alla tua famiglia
. Questo detto è l'unico che credo non seguirò almeno non con i miei fratelli. Ti ringrazio per tutto e sarai sempre una persona speciale per me. Concludo dicendo che una volta mi domandasti quale sarebbe stata la mia reazione in un ipotetica rissa dove qualcuno voleva farti del male, all'epoca non ti diedi risposta perché non so il mio corpo come potrebbe reagire ad una simile situazione, continuo a non avere quella risposta ma lasciami dire solo una cosa, se dovessi essere lì con te non indietreggerei.\\

La famiglia Pellecchia è stata la prima che mi ha trattato come un figlio, vorrei ringraziare Marianna e anche Antonio per avermi fatto sentire in una seconda casa, in una seconda famiglia. Vedete ero sempre scettico sul fatto di non voler disturbare all'inizio quando ricevevo inviti a venire a casa perché sapevo che potevo dare fastidio, invece non è mai successo anzi era proprio il contrario! Mai all'epoca mi sarei ricreduto su una simile cosa. Questa famiglia ha una cosa che credo leghi bene con la mia e che in un certo senso mi ha permesso di capire Quanto io sia stato fortunato, non solo ho trovato un fratello quì ma anche una donna che come mia madre ne ha passate tante e non si arrende alla vita, anche se quest'ultima continua e continua a colpire, lei per il bene dei suoi figli continua a combattere. Per non dimenticare le sorelle, Sara in grado di essere credo a questo punto una chef mancata dato che cucina da dio e Roberta, ricordo che ancora dovevi nascere, invece adesso sei cresciuta parecchio, comunque apro e chiudo parentesi si vede che hai preso da tuo fratello a cacciarti nei guai. Insomma una seconda famiglia che ho trovato!\\

Passiamo ora ai Bravaccini non ultimi perchè meno importanti eh! ma non posso mettere il testo uno di fianco all'altro chiedo venia, parto da david.\\

Parto da un mio secondo fratello non consanguineo, David Bravaccini.\\

David ormai possiamo dire che 5 anni sono passati da quando in quel fatidico giorno di estate ci siamo incontrati a Miliscola, come ti ho sempre detto non mi hai fatto una bella impressione all'inizio, oddio il classico palestrato e belloccio che si vantava di qualsiasi cosa, almeno questo fino a 5 secondi prima che cominciassimo a parlare e a conoscerci, non mi ero mai sbagliato così tanto su una persona, incredibile che eri totalmente diverso da ciò che vedevo. Cosa ancora più strana è la velocità con cui siamo diventati amici all'inizio e poi fratelli, insomma quasi istantanea la cosa. Anche con te ne ho passate, eccome se ne ho passate, e tutti e tre meglio che non ne parliamo proprio, si è innescata subito una miccia che continua a bruciare, anche tu sei stato al mio fianco quando comunque dovetti affrontare determinate cose, sei stato il primo con cui ho confidato e ho parlato del fatidico momento di buio (dato che Nikol me lo aveva proprio visto vivere), senza mai dirmi nulla di offensivo e senza mai dilungarti in critiche, hai solo ascoltato quello che avevo da dire e mi è bastato per capire che persona sei. Metti sempre d'avanti le persone che vuoi bene, un leone che sbrana chiunque voglia far loro del male, sicuro di te a livelli che veramente sembra che non ti importi niente delle conseguenze di quello che fai ma forse questa è una delle tue peculiarità che ammiro molto, anche se a volte richiama più guai che altro ma tanto per la maggior parte da quando ci siamo conosciuti in quei guai stavo pure io, insomma non c'è 2 senza 3 no? Non vorrei dilungarmi troppo dato che ciò che hai fatto per me non lo dimenticherò mai e tu dirai perché cosa ho fatto? come cosa, mi hai fatto conoscere altre persone fantastiche, mi hai incluso sempre in qualsiasi cosa che sia un uscita di famiglia o un uscita con gli altri del gruppo o altre persone. Il primo viaggio ufficiale diciamo per così dire che abbiamo fatto tutti e tre ad Amsterdam, quella si che è un bellissimo ricordo che mai dimenticherò. Insomma sono sempre stato lì non perché mi auto-invitavo o chiedevo di partecipare, ma, perché già c'era una persona che metteva in conto che io dovevo esserci. Il tatuaggio che ci siamo fatti sancisce questa solidarietà fraterna solo simbolicamente, siamo noi che lo rendiamo vivo con i valori che ci siamo imposti o che ci hanno insegnato, Sapendo che possiamo fidarci l’uno dell’altro. Insomma vale ciò che ho detto a Nikol, ossia che non so se riuscirò ad essere colui che vi aspettate che io sia, ma, non vi volterò mai le spalle.\\

Nella famiglia Bravaccini ho incontrato persone con un'umanità speciale, anche loro mi hanno trattato come un figlio senza mai farmi pesare le innumerevoli volte che mi hanno accolto a casa loro anzi facendomi sentire proprio a casa! Ho capito una cosa fondamentale che ci unisce ossia l'amore delle nostre madri che provano per noi, insomma Marianna e Monica per come affrontano la vita mi ricordano mia madre, anche loro hanno dovuto patire la scelta di dover affrontare la maggior parte delle cose da sole, forse l'unica differenza è che almeno David ha ancora un padre presente, però non mi dilungo sulla questione dato che voglio ringraziare tutte le persone per ciò che hanno fatto per me. I miei più sentiti ringraziamenti vanno in primis a Monica e Gaetano, i genitori di David che mi hanno fatto sentire realmente parte della loro famiglia senza dover essere per forza un loro parente, A Deborah e a Francesco che per un periodo sono stati dei veri e propri "Masti" che mi hanno insegnato a fare tante cose nel loro laboratorio e non solo per questo ma anche per i motivi citati sopra, ai nonni che anche loro stavano aspettando questo fatidico giorno e che finalmente è arrivato, insomma non riuscirò mai a descrivere la mia gratitudine per questo continuerò ad esservi per sempre grati.\\


Ringrazio tutti i miei amici, colleghi universitari, le persone fantastiche che ho incontrato a lavoro, come il mio gruppo di TA (Francesca Di Vincenzo "senpai", Francesco Giamminelli, Pasquale Cioffi, Martin Cioffi "il mentore di mooney",  Alessandro Leopardi etc..)che mi stanno facendo amare questo lavoro e quest'azienda, ogni giorno muoio dalle risate a furia di parlare con loro. Ringrazio i miei amici del corso che poi sono stati anche loro ammessi in Accenture, e ringrazio anche le altre persone incontrate lì (Marco Fiorillo, Pasquale Giusta, Francesco Scala, Raffaella Della Gatta, Mattia Gambardella "Fratm", Ciro Montagna, Emilia Cammisa, Annalisa Desiato etc...) non metto tutti altrimenti non finisco più e tutti i miei amici più cari che non sto quì a scrivere altrimenti ribadisco non finisco più.\\

ai miei colleghi universitari come Giovanni Belluomo che abbiamo veramente rasentato la pazzia a furia di studiare, a Davide Russo e Salvatore Pappagallo che ne abbiamo combinate delle grosse all'inizio della nostra carriera universitaria, a DAvide De Angelis che mi ha aiutato a non farmi venire un'infarto per prenotare la seduta di laurea e al vecchio gruppo che con tutto che non ci sentiamo più è doveroso ringraziare.


 ad Emmanuele Perrella mio collega e amico la cui caratteristica oltre ad essere una statua scolpita nel marmo è di un'umanità incredibile, in 1 anno che ci siamo conosciuti mi ha permesso di condividere tante di quelle risate e di avere una persona più che disponibile nel condividere sapere e idee dalle più serie alle più disparate, un vero esempio da seguire o comunque da prendere spunto date che grazie a lui sono riuscito anche a prendermi un vestito decente! Scherzi a parte veramente una persona dalla bontà più unica che rara, pure sul lavoro sono stato molto fortunato.\\

a Riccardo Storti colui che mi ha permesso di entrare a far parte di questa grande azienda, che ha creduto in me fino alla fine e che è sempre stato disponibile a qualsiasi chiarimento e a qualsiasi dubbio che la mia mente potesse partorire, sei stato colui che mi ha dato un posto indeterminato non so come ringraziarti!\\

vorrei ringraziare anche Nicole Novi, una donna che ne ha passate tante, che veramente dio solo sa come diavolo ha fatto a rialzarsi da qualsiasi torto che ha subito. Mi hai dato tanti consigli preziosi anche per la tesi e per questo fine percorso, ho potuto scambiare idee e pareri con una persona che ha già affrontato ciò che tocca affrontare me e questa volta li ho seguiti i consigli. Insomma spero veramente che la vita ti sorrida e che tu riesca a trovare la tua felicità, forse una delle persone più buone che abbia mai incontrato e non è poco. Ti ringrazio anche perché in questo ultimo anno alle prese col l'esame più difficile che io abbia mai affrontato sei stata capace di farmi ragionare e di farmi vedere le cose da una prospettiva diversa, non so come sia possibile ma ci sei riuscita e per questo non so come ringraziarti se non con un generico ma non banale Grazie!\\

I miei più sentiti ringraziamenti vanno alla relatrice e il correlatore corrispettivamente Livia Marcellino e Pasquale De Luca per avermi seguito in questo step finale, con un'umanità indescrivibile e una disponibilità che mi ha lasciato senza parole. Provenendo da un periodo di stagnazione del mio percorso formativo, non credevo di ritrovare docenti di riferimento così accorti, capaci non solo di guidarmi dal punto di vista accademico, ma anche di restituirmi fiducia nelle mie capacità e nel valore del mio percorso. Il loro supporto costante, unito a una sincera attenzione verso la mia crescita personale e professionale, ha rappresentato uno stimolo fondamentale per portare a termine questo lavoro con rinnovata motivazione ed entusiasmo.
